\chapter{Potenzmengen}

\section{Teilmengen als Elemente}

\begin{definition}[Potenzmenge]
Die \emph{Potenzmenge} von $A$ ist die eindeutig bestimmte Menge
$\mathcal P(A)$ mit
\[
 X\in\mathcal P(A)\Longleftrightarrow X\subseteq A.
\]
\end{definition}

\begin{theorem}[Grundlegende Elemente der Potenzmenge]
\[
 \varnothing\in\mathcal P(A),\qquad A\in\mathcal P(A),\qquad
 \mathcal P(A)\ne\varnothing.
\]
\end{theorem}

\begin{theorem}[Inklusion und Gleichheit von Potenzmengen]
\label{b03guide:thm:powercomparison}
\[
\begin{aligned}
 A\subseteq B&\Longleftrightarrow
                     \mathcal P(A)\subseteq\mathcal P(B),\\
 A=B&\Longleftrightarrow\mathcal P(A)=\mathcal P(B),\\
 A\subsetneq B&\Longleftrightarrow
                     \mathcal P(A)\subsetneq\mathcal P(B).
\end{aligned}
\]
Hier geht es um tatsächliche Inklusion und Gleichheit von Mengen.
\cite{hammack}
\end{theorem}

\begin{theorem}[Vereinigung und Durchschnitt aller Teilmengen]
\[
 \bigcup\mathcal P(A)=A,\qquad \bigcap\mathcal P(A)=\varnothing.
\]
Der Durchschnitt ist stets definiert, da $\mathcal P(A)$ nichtleer ist.
\end{theorem}

\begin{theorem}[Jede Familie besitzt eine gemeinsame umfassende Menge]
\[
 \mathcal F\subseteq\mathcal P(U)
       \Longleftrightarrow \bigcup\mathcal F\subseteq U.
\]
Insbesondere gilt
$\mathcal F\subseteq\mathcal P(\bigcup\mathcal F)$.
\end{theorem}

\section{Potenzmengen und Mengenoperationen}

\begin{theorem}[Potenzmenge eines Schnitts]
\[
 \mathcal P(A\cap B)=\mathcal P(A)\cap\mathcal P(B).
\]
Allgemeiner gilt für $\mathcal F\ne\varnothing$
\[
 \mathcal P(\bigcap\mathcal F)
       =\bigcap\{\mathcal P(A)\mid A\in\mathcal F\}.
\]
\end{theorem}

\begin{theorem}[Potenzmenge einer Vereinigung]
\[
 \mathcal P(A)\cup\mathcal P(B)\subseteq\mathcal P(A\cup B).
\]
Es gilt genau dann Gleichheit, wenn $A\subseteq B$ oder
$B\subseteq A$. Ohne diese Bedingung gilt stets die Darstellung
\[
 \mathcal P(A\cup B)
   =\{X\cup Y\mid X\in\mathcal P(A),\ Y\in\mathcal P(B)\}.
\]
Für jede Familie $\mathcal F$ gilt
\[
 \bigcup\{\mathcal P(A)\mid A\in\mathcal F\}
                      \subseteq\mathcal P(\bigcup\mathcal F).
\]
In dieser Inklusion gilt genau dann Gleichheit, wenn
$\bigcup\mathcal F\in\mathcal F$; dies schließt den Fall der
leeren Familie in die Aussage ein.
\end{theorem}

\begin{theorem}[Alle Teilmengen einer Familienvereinigung]
Für jede Familie $\mathcal F$ sei
$\mathcal D=\bigcup\{\mathcal P(A)\mid A\in\mathcal F\}$.
Dann gilt
\[
 \mathcal P(\bigcup\mathcal F)
   =\{\bigcup\mathcal G\mid\mathcal G\in\mathcal P(\mathcal D)\}.
\]
Rechts werden sämtliche Teilfamilien von $\mathcal D$ berücksichtigt;
die leere Familie $\mathcal F$ ist eingeschlossen.
\cite{koepke,hammack}
\end{theorem}

\begin{theorem}[Potenzmenge einer Differenz]
\[
 \mathcal P(A\setminus B)
       =\{X\in\mathcal P(A)\mid X\cap B=\varnothing\}.
\]
Für beliebige Mengen $A,B$ gilt ferner
\[
 \{A\setminus X\mid X\in\mathcal P(B)\}
       =\{Y\in\mathcal P(A)\mid A\setminus B\subseteq Y\}.
\]
\end{theorem}

\enlargethispage{\baselineskip}
\begin{theorem}[Potenzmenge des relativen Durchschnitts]
Für jede Menge $U$ und jede Familie $\mathcal F$ gilt
\[
 \mathcal P(\bigcap\nolimits_U\mathcal F)
    =\bigcap\nolimits_{\mathcal P(U)}
                         \{\mathcal P(A)\mid A\in\mathcal F\}.
\]
Dies gilt auch bei leerer Familie.
\end{theorem}

